%%%%%%%%%%  scan idx 614  |  folio 601  %%%%%%%%%%
\((X_i)_\eta\cap(X_j)_\eta=\varnothing\). Soient \(Y_k=\bigcup_{i\in I(k)}X_i\) les composantes connexes de \(X\)~; alors \(f\) se décompose, \(f=\coprod_k f_k:\coprod_k(\coprod_{i\in I(k)}X_i)\longrightarrow Y_k\), et l'assertion en résulte. Enfin, une récurrence immédiate nous ramène au cas d'un morphisme \(X_1\amalg X_2\to X\). \(\sgacircled{1}\) et \(\sgacircled{3}\) résultent alors, quitte à restreindre \(S\) d'après (EGA III 7.8.6), de la

\textbf{Proposition 2.3.} \uline{Soient} \(S\) \uline{un schéma,} \(g:Y\to S\) \uline{un morphisme propre et} \(Y_1,Y_2\) \uline{deux sous-schémas fermés définis par les idéaux} \(I_1,I_2\)\uline{. On suppose que}
\begin{enumerate}[label=\alph*)]
\item \(Y,Y_1,Y_2,Y_1\cap Y_2\) \uline{sont fidèlement plats et de présentation finie sur} \(S\).
\item \(g_*{O}_Y={O}_S\) et \((g_i)_*{O}_{Y_i}={O}_S\) \uline{universellement (où} \(g_i=g|Y_i\), \(i=1,2\)).
\item \(I_1\cap I_2=0\).
\end{enumerate}
\uline{Alors} \(\underline{\Pic}_{Y/S}\longrightarrow \underline{\Pic}_{(Y_1\amalg Y_2)/S}\) \uline{est représentable par un morphisme affine de présentation finie.}

En effet notons d'abord que les hypothèses sont stables pour tout changement de base. Pour (c), celà résulte des deux suites exactes~:
\[
  0\longrightarrow I_1+I_2\longrightarrow{O}_Y
  \longrightarrow{O}_{Y_1\cap Y_2}\longrightarrow0
\]
et
\[
  0\longrightarrow I_1\cap I_2\longrightarrow I_1\oplus I_2
  \longrightarrow I_1+I_2\longrightarrow0~;
\]
la première montre que \(I_1+I_2\) est plat sur \(S\)~; ensuite la seconde que \(I_1\cap I_2\) est stable par changement de base.

L'assertion est locale pour la topologie f.p.p.f. En faisant le changement de base \(Y_1\cap Y_2\to S\), on peut donc supposer que \(Y_1\cap Y_2/S\) possède une section \(\alpha\), qui définit des sections \(\alpha_i\) de \(Y_i\) et \(\beta\) de \(Y\).

Après un changement de base convenable, on est ramené à voir que tout point \(\mu:S\to\underline{\Pic}_{(Y_1\amalg Y_2)/S}\) a comme image réciproque un sous-foncteur de
